\section{Skorohod and Stratonovich wave equations}\label{S:literature-context}

The meaning of the noise product is part of the model. The same formal wave
operator and covariance do not specify the same equation under the Skorohod
and Stratonovich interpretations. We review both interpretations below,
separating existence results, moment-growth bounds and exact asymptotics.
All comparisons retain the hypotheses of the cited papers.

\subsection{The two products and their consequences}

The Skorohod integral is the divergence operator of Malliavin calculus. For the
linear equation with deterministic initial data, its iterated integrals form
orthogonal homogeneous Wiener chaoses. In the convergent expansion, every
positive-order term has mean zero, so the mean solution is the deterministic
free wave. Orthogonality reduces its second moment to a sum of squared chaos
norms. For adapted integrands and white-in-time noise, this interpretation
agrees with the It\^o integral in the corresponding integration framework;
see Walsh~\cite{walsh:86:introduction} and Dalang~\cite{dalang:99:extending}.
Temporal correlation removes the martingale structure, even when the
Skorohod integral remains available.

The Stratonovich construction considered here takes ordinary products of
regularized Gaussian noise and then removes the regularization. Wick's
formula retains contractions within an iterated term as well as contractions
between different terms. These self-pairings change the mean; terms of unequal
iteration orders can contribute to the second moment. This is why the
Skorohod second-moment isometry does not directly solve the present problem.
It does not mean that a Stratonovich solution has no orthogonal chaos
expansion: trace terms can be regrouped into Wiener chaoses, as in
Balan~\cite{balan:22:stratonovich}. The iteration order and the final Wiener
chaos order are different indices.

For regular noise, ordinary multiplication connects the equation to
deterministic wave evolution in a random potential. Stratonovich models are
therefore natural for specified smooth-noise approximations. For singular noise, convergence, the approximation scheme and any
required counterterms must be proved. There is no universal conversion formula
that identifies the two wave equations. In particular, an adapted
white-in-time Skorohod solution being an It\^o solution does not identify it
with an unspecified smooth-noise limit. The three-dimensional approximation
analysis of Delgado-Vences--Sanz-Sol\'e
\cite{delgado-vences.sanz-sole:14:approximation} shows why
regularization requires care.

\begin{figure}
\centering
\begin{tikzpicture}[x=1cm,y=1cm]
  \node[ham/box,minimum width=5.2cm] (hamSkoroProduct) at (0,0)
  {Skorohod / Wick product\\$u\diamond\dot W$};
  \node[ham/box,minimum width=5.2cm] (hamStratoProduct) at (6.5,0)
  {Ordinary Stratonovich construction\\regularize first, then multiply};
  \node[ham/box,minimum width=5.2cm] (hamSkoroIntegral) at (0,-1.8)
  {Divergence stochastic integral\\and Wiener-chaos expansion};
  \node[ham/box,minimum width=5.2cm] (hamStratoLimit) at (6.5,-1.8)
  {Classical random wave equations\\followed by the specified limit};
  \draw[ham/arrow] (hamSkoroProduct)--(hamSkoroIntegral);
  \draw[ham/arrow] (hamStratoProduct)--(hamStratoLimit);
  \node[ham/note] at (0,-3.6)
  {Adapted white-in-time integrands:\\
  Skorohod agrees with It\^o\\
  under the integration hypotheses.};
  \node[ham/note] at (6.5,-3.6)
  {Singular noise requires\\
  a specified approximation scheme\\
  and a convergence theorem.\\
  Correction terms may be needed.};
  \draw[hamblue!35] (-2.8,-5.0)--(9.3,-5.0);
  \node[ham/box,minimum width=10.9cm] at (3.25,-6.0)
  {The same formal wave operator and noise covariance\\
  do not identify the solutions or their moment constants.};
\end{tikzpicture}
\caption[Two meanings of the noise product]{The stochastic product is part of the model. The arrows run only
within each construction; there is no conversion arrow between the columns.
For the linear Skorohod equation with deterministic data, iteration order
agrees with orthogonal chaos order. Ordinary Stratonovich iteration retains
self-contractions and unequal-order moment terms, although the final
solution can still be regrouped into orthogonal Wiener chaoses. See
Balan~\cite{balan:22:stratonovich} for a comparison of the expansions.
An adapted It\^o identification does not identify an unspecified
smooth-noise approximation.}
\label{Fig:literature-interpretations}
\end{figure}


\subsection{White-in-time noise}

White-in-time wave equations provide a useful starting point because adapted
stochastic integration is available. Conus--Dalang
\cite{conus.dalang:08:non-linear} develop constructions in higher dimensions
under their noise and initial-data assumptions. Dalang--Mueller--Tribe
\cite{dalang.mueller.ea:08:feynman-kac-type} obtain Feynman--Kac-type
representations using Poisson processes and random flights. Thus, the absence
of the usual Brownian heat-equation representation does not imply the absence
of all probabilistic moment representations for wave equations.

Chen--Dalang~\cite{chen.dalang:15:moment}*{Theorem~2.3 and Corollary~1.1}
derive moment bounds and explicit second-moment formulas for the
one-dimensional space-time white-noise equation, including rough initial
data. Chen--Guo--Song~\cite{chen.guo.ea:24:moments}*{Theorem~1.1} study a
fractional SPDE family driven by space-time white noise. Their second moments
and second Lyapunov exponents are explicit. Their wave specialization gives
matching upper and lower $p^{3/2}$ moment-growth orders under the stated data
conditions~\cite{chen.guo.ea:24:moments}*{Proposition~5.10; (1.10)--(1.11),
pp.~4257--4258}. The higher-moment constants are not optimal, and these bounds
do not assert an exact higher-moment leading constant or a general
higher-moment limit.

Balan--Song~\cite{balan.song:19:second}*{Theorems~1.1 and 1.3} relate the
second-moment Laplace transforms of hyperbolic and parabolic Anderson models.
Their general result identifies an upper second Lyapunov exponent; their
Riesz specialization proves an actual limit. These conclusions should not be
interchanged. In three dimensions with a different spatial covariance class,
Dalang--Mueller~\cite{dalang.mueller:09:intermittency} establish upper and
lower moment-growth bounds of order $p^{4/3}$, with the lower bound at even
orders and with separate hypotheses on the two bounds. The power of the moment
order therefore depends on the spatial covariance, not only on the word
``wave''.

Rough spatial noise is another distinct direction. Balan--Jolis--Quer-Sardanyons
\cite{balan.jolis.ea:17:intermittency} treat one-dimensional spatial fractional
noise with Hurst index between $1/4$ and $1/2$; Liu--Hu--Wang
\cite{liu.hu.ea:22:nonlinear} develop nonlinear well-posedness in this range
under their coefficient assumptions. Such distributional spatial covariances
are outside the nonnegative spatial-measure class of the present paper. Their
results do not extend the scope of Hypothesis~\ref{H:homog} by analogy.

\subsection{Time-independent noise}

For the static Skorohod equation with initial data $(1,0)$, Balan--Le Chen--Xia
Chen~\cite{balan.chen.ea:22:exact}*{Theorem~1.6} obtain exact large-time and
high-moment asymptotics with a variational constant. Their homogeneous spatial
covariances cover dimensions at most three; their global spatial white-noise
case covers dimensions one and two. Three-dimensional white noise belongs to
a separate finite-time discussion, not this global large-time assertion.
Chen--Eisenberg~\cite{chen.eisenberg:23:interpolating}*{Theorems~1.6--1.7}
study fractional differential operators with static Skorohod noise. Their time
order lies strictly below two. Fractional time differentiation is distinct
from temporal correlation of the noise and does not silently include the
ordinary wave endpoint.

For ordinary static Stratonovich products, Balan
\cite{balan:22:stratonovich} constructs solutions in dimensions one and two
under a spectral condition stronger than Dalang's condition, and gives both
a probabilistic representation and a comparison of chaos expansions.
Xia Chen--Hu~\cite{chen.hu:25:hyperbolic} treat dimensions at most three under
Dalang's condition and establish homogeneous moment asymptotics at positive
integer orders. Their published variational expressions require the
normalization qualifications recorded in this manuscript; we do not import
a numerical white-noise constant from them.

There is a concrete difference in admissible roughness for these unrenormalized
static constructions. For the common flat data and homogeneous function
covariances, the Skorohod range in the preceding exact-asymptotics result
allows spatial homogeneity exponent below three, subject to its dimension
restriction, whereas the ordinary Stratonovich range requires exponent below
two. This comparison concerns those solution classes. Renormalized
Anderson-Hamiltonian wave equations, such as those in
Gubinelli--Ugurcan--Zachhuber~\cite{gubinelli.ugurcan.ea:20:semilinear}, are a
related but different construction; their existence does not identify an
ordinary Stratonovich product with a Skorohod product.

\subsection{Fractional temporal correlation}

On the Skorohod side, Balan--Song~\cite{balan.song:17:hyperbolic} develop the
space-time homogeneous Gaussian-noise construction. Its higher-dimensional
statements have additional spectral hypotheses, including absolute continuity
of the spectral measure in the relevant results. Xia Chen--Deya--Song--Tindel
\cite{chen.deya.ea:25:solving} sharpen the spectral solvability analysis. For
homogeneous spatial covariance and temporal singularity exponent $\alpha_0$,
the resulting noise threshold is $\alpha + \alpha_0 < 3$, under the paper's
stated data conditions. The sufficiency statement in dimensions two and three
imposes Fourier integrability on the initial displacement; the necessity
statement uses a positive lower bound on the free wave. These are different
hypotheses and directions, not an unconditional common-data equivalence with
the present construction. The static endpoint also needs separate treatment.

Balan--Conus~\cite{balan.conus:16:intermittency} give moment bounds for
fractional-in-time noise under their spatial covariance conditions, including
a second-moment lower bound in dimensions at most three at the same time-growth
order as the upper bound for positive constant initial displacement and
nonnegative constant initial velocity. Hu--Wang~\cite{hu.wang:24:matching}*{Theorem~3.7}
obtain matching higher-moment growth orders under their covariance and
Green-kernel hypotheses. Neither conclusion by itself identifies a sharp
leading constant. Space-time white noise is not an automatic specialization
of their functional spatial lower-bound hypotheses.

On the Stratonovich side, Xia Chen--Deya--Song--Tindel
\cite{chen.deya.ea:23:hyperbolic} construct solutions in dimensions one and two
using Strichartz estimates and weighted function spaces. Their sufficient
regularity ranges and initial-data assumptions differ from the full
homogeneous range considered here. Xia Chen's current author version
\cite{chen:author2026:time-dependency} treats fractional temporal correlation
in dimensions one through three, with homogeneous condition
$\alpha + \alpha_0 < 2$. It proves two-sided first-moment asymptotic bounds and
a coarse quadratic-exponential second-moment upper bound. The stronger exact
first-moment statement of arXiv:2510.01412v1 is not the statement of this
edition; the version convention is specified in Section~\ref{S:known}.

\begin{figure}
  \centering
  \begin{tikzpicture}[x=1cm,y=1cm,
    lit/text/.style={anchor=north west,align=flush left,text width=5.35cm,
      font=\small,inner sep=0pt,
      execute at begin node={\hyphenpenalty=10000\relax}},
    lit/tag/.style={anchor=west,font=\small\bfseries,inner sep=0pt},
    lit/regime/.style={align=center,font=\small,inner sep=0pt}]
    % The two colored columns classify products, not solution inclusions.
    \fill[hamblue!3] (1.65,.35) rectangle (7.35,-11.75);
    \fill[hamteal!3] (7.75,.35) rectangle (13.45,-11.75);
    \fill[hamblue] (1.65,.35) rectangle (7.35,1.15);
    \fill[hamteal] (7.75,.35) rectangle (13.45,1.15);
    \node[text=white,font=\bfseries] at (4.5,.75) {Skorohod};
    \node[text=white,font=\bfseries] at (10.6,.75) {Ordinary Stratonovich};
    \node[lit/regime,text=hamblue!80!black] at (.55,.75)
      {Temporal\\covariance};

    % White time: the spike represents a delta covariance, not a noise path.
    \node[lit/regime] at (.55,-.25) {White\\in time};
    \draw[gray] (-.05,-1.4)--(1.15,-1.4);
    \draw[-{Latex[length=1.6mm]},thick,hamgold]
      (.55,-1.4)--(.55,-.85);
    \node[lit/regime] at (.55,-1.8) {$\delta_0$};
    \node[lit/tag,text=hamblue] at (1.9,-.1) {Exact second moments};
    \node[lit/text] at (1.9,-.45)
      {Chen--Dalang~\cite{chen.dalang:15:moment}\\
      Space-time white noise, $d = 1$.\\[4pt]
      Chen--Guo--Song~\cite{chen.guo.ea:24:moments}\\
      Exact $p = 2$; wave higher-$p$\\
      growth orders $p^{3/2}$.};
    \node[lit/tag,text=hamteal] at (8,-.1) {Approximation matters};
    \node[lit/text] at (8,-.45)
      {No source claim is listed here.\\[4pt]
      Specify smoothing and convergence.\\[4pt]
      No It\^o equality or openness\\
      is inferred.};

    \draw[hamblue!18] (1.9,-3.2)--(7.1,-3.2);
    \draw[hamteal!22] (8,-3.2)--(13.2,-3.2);
    % Static time: one constant covariance profile.
    \node[lit/regime] at (.55,-3.65) {Static\\noise};
    \draw[gray] (-.05,-4.8)--(1.15,-4.8);
    \draw[thick,hamgold] (-.05,-4.25)--(1.15,-4.25);
    \node[lit/regime] at (.55,-5.2) {$1$};
    \node[lit/tag,text=hamblue] at (1.9,-3.7) {Exact rates};
    \node[lit/text] at (1.9,-4.05)
      {Balan--Le Chen--Xia Chen\\
      \cite{balan.chen.ea:22:exact}: $p \ge 2$.\\
      Homogeneous: $d \le 3$;\\
      spatial white: $d \le 2$.\\[4pt]
      Chen--Eisenberg~\cite{chen.eisenberg:23:interpolating}\\
      Exact rates; static noise;\\
      fractional time order $< 2$.};
    \node[lit/tag,text=hamteal] at (8,-3.7) {Construction / exact rates};
    \node[lit/text] at (8,-4.05)
      {Balan~\cite{balan:22:stratonovich}: $d = 1,2$ construction.\\
      Stronger spectral condition.\\[4pt]
      Xia Chen--Hu~\cite{chen.hu:25:hyperbolic}\\
      Exact rates, integer $p \ge 1$;\\
      Dalang and homogeneity hypotheses.};

    \draw[hamblue!18] (1.9,-7.7)--(7.1,-7.7);
    \draw[hamteal!22] (8,-7.7)--(13.2,-7.7);
    % Fractional time: schematic singular, decreasing covariance.
    \node[lit/regime] at (.55,-8.45) {Fractional\\in time};
    \draw[gray] (-.05,-9.65)--(1.15,-9.65);
    \draw[thick,hamgold]
      (.05,-8.85) .. controls (.18,-9.35) and (.65,-9.45) .. (1.15,-9.52);
    \node[lit/regime] at (.55,-10.1) {$|t|^{-\alpha_0}$};
    \node[lit/tag,text=hamblue] at (1.9,-8.1) {Bounds / construction};
    \node[lit/text] at (1.9,-8.45)
      {Balan--Conus; Hu--Wang\\
      \cites{balan.conus:16:intermittency,hu.wang:24:matching}\\
      Growth orders; no sharp constant.\\[4pt]
      Chen--Deya--Song--Tindel\\
      \cite{chen.deya.ea:25:solving}: construction;\\
      data hypotheses depend on $d$.};
    \node[lit/tag,text=hamteal] at (8,-8.1) {First-moment bounds};
    \node[lit/text] at (8,-8.45)
      {Xia Chen, pinned author version\\
      \cite{chen:author2026:time-dependency}\\
      $d = 1,2,3$, $\alpha + \alpha_0 < 2$.\\[4pt]
      Two-sided first-moment bounds;\\
      the exact limit is predicted.};

    % A separate manuscript contribution, not a transferred literature claim.
    \fill[hamgold!9,rounded corners=2pt]
      (1.65,-12.1) rectangle (13.45,-13.65);
    \draw[hamgold,line width=2pt] (1.65,-12.1)--(1.65,-13.65);
    \node[lit/tag,text=hamgold!85!black] at (1.95,-12.45)
      {Present manuscript \quad Fractional-time Stratonovich, $p = 2$};
    \node[anchor=north west,align=left,text width=10.9cm,
      font=\small,inner sep=0pt] at (1.95,-12.75)
      {Intrinsic variational formula for the second-moment limit.\\
      Written proof; independent review remains pending.};
  \end{tikzpicture}
  \caption[Source-qualified literature comparison]{Selected results, organized
    by temporal covariance and stochastic product. The small profiles are
    schematic covariance sketches, not sample noise paths or limiting arrows.
    Construction, growth orders and exact rates are different conclusions.
    Every source retains its spatial covariance, dimension and data hypotheses;
    the Chen--Guo--Song higher-moment constants are nonoptimal, and the
    higher-dimensional Skorohod construction retains its data qualifications.
    Chen--Hu supplies no numerical constant here. Xia Chen's pinned author
    version gives first-moment bounds, unlike the stronger arXiv-v1 wording.
    The highlighted manuscript claim has a separate review status. The columns
    assert neither inclusions nor a universal conversion between products.}
  \label{Fig:literature-results-map}
\end{figure}


\subsection{Advantages, limitations and the present contribution}

For second-moment analysis, the main advantage of the Skorohod formulation is
orthogonality: it removes unequal-chaos-order cross terms and makes isometries
and hypercontractivity directly available. In some static covariance classes
it also permits rougher noise than ordinary Stratonovich multiplication. The
price is that Wick-ordered multiplication specifies a different stochastic model;
it cannot generally be substituted for an ordinary smooth random potential.
With temporally correlated noise, it does not restore adapted martingale
methods.

The Stratonovich formulation retains the interactions generated by ordinary
multiplication and can connect to deterministic wave and Gaussian variational
methods. Its cost is control of traces, self-pairings and unequal iteration
orders, together with potentially stronger conditions for an unrenormalized
limit. Neither interpretation supplies a general wave maximum principle.
Both admit chaos methods and, in suitable settings, probabilistic
representations. The appropriate choice is determined by the model and the
specified approximation, rather than by a universal preference for one
calculus.

The present contribution concerns the temporally correlated Stratonovich
problem. Theorem~\ref{T:full-upper} gives the full homogeneous-class upper bound
at the scale $t^\kappa$. Theorem~\ref{T:intrinsic-second-moment} identifies the
second-moment limit through nonnegative compact space-time controls, retaining
the original covariance amplitude, unit kinetic coefficient and global
half-variance cost. This intrinsic expression is not a numerical closed form,
a sharp all-moment formula or an identification with the source's
first-moment variational constants. These are the claims of the written
proofs; the independent review status is recorded separately from this
literature comparison.

Related work in progress with Raluca Balan and Jinxin Wang treats three
Skorohod projects: noise colored in time and space in one dimension; colored
noise in two dimensions; and white-in-time noise in dimensions one, two and
three. We mention these scopes to locate the present work within the broader
program. No unpublished result from them is used here, and no equality of
their constants with the present Stratonovich constant is asserted.
